% Generated draft: requires mathematical and source review.
\documentclass[11pt]{article}
\usepackage[margin=25mm]{geometry}
\usepackage{fontspec}
\setmainfont{Latin Modern Roman}
\usepackage{amsmath,amssymb,mathtools,unicode-math}
\setmathfont{Latin Modern Math}
\usepackage{enumitem,needspace,xurl,hyperref}
\setlength{\parindent}{0pt}
\setlength{\parskip}{5pt}
\setlength{\emergencystretch}{3em}
\newcommand{\N}{\mathbb{N}}
\newcommand{\Z}{\mathbb{Z}}
\newcommand{\Q}{\mathbb{Q}}
\newcommand{\R}{\mathbb{R}}
\newcommand{\C}{\mathbb{C}}
\newcommand{\F}{\mathbb{F}}
\newcommand{\A}{\mathbb{A}}
\newcommand{\PP}{\mathbb{P}}
\newcommand{\E}{\mathbb{E}}
\newcommand{\eps}{\varepsilon}
\newcommand{\dd}{\,\mathrm{d}}
\begin{document}
\section*{UnsolvedMath ID: 30006436}
\textbf{Source problem code:} catalogue entry
\subsection{Definitions and mathematical statement}
Let \(\mathfrak H=\{x+iy\in\C:y>0\}\). For reference, a standard weight-\(k\) hyperbolic Laplacian is
\[
\Delta_k=-y^2\left(\frac{\partial^2}{\partial x^2}+\frac{\partial^2}{\partial y^2}\right)+iky\left(\frac{\partial}{\partial x}+i\frac{\partial}{\partial y}\right).
\]
Once a modular group, weight, automorphy law, and cusp-growth conditions have been fixed, the zero-eigenvalue polyharmonic condition of order at most \(r\) on a smooth automorphic function \(F\) is
\[
\Delta_k^{\,r}F=0,\qquad r\in\Z_{\geq 1}.
\]
The harmonic condition is the case \(r=1\). This fixes terminology for the differential equation only: the supplied record does not specify the automorphy data, the growth conditions meant by ``weak,'' or the intended conventions for polyharmonic order. No equality between polyharmonic order and mock-modular depth is assumed.

For a real Lie group \(G\) with maximal compact subgroup \(K\), write \(\mathfrak g\) for its complexified Lie algebra. A \((\mathfrak g,K)\)-module carries compatible Lie-algebra and locally finite \(K\)-actions. Such a module has finite length if it has a finite filtration by submodules whose successive quotients are irreducible. The supplied record does not identify the precise ambient group, module categories, or meaning of ``well-behaved'' representation theory.

\textbf{Research problem.} Construct and specify a connection between higher-depth mixed mock modular forms, their iterated-integral descriptions and modular completions, and polyharmonic weak Maass forms. The construction should state its admissible groups, weights, depths, automorphy data, cusp conditions, and any relation between depth and polyharmonic order. Use the connection to explain the requested passage from finite-length representation theory associated with the modular group to a well-behaved representation theory for the ambient real Lie group. Here the requested passage is a research objective; the source supplies neither a functor nor a precise categorical property that would constitute its resolution.

As the initial case, determine whether the iterated Eisenstein integrals arising from the motivic viewpoint correspond directly to polyharmonic Eisenstein series arising from the analytic viewpoint, or whether cusp forms or weak forms enter as correction terms or obstructions. The source explicitly leaves this alternative unresolved. It specifies no integral lengths, regularizations, normalizations, Eisenstein parameters, or permitted correction spaces, so it does not determine a particular identity or a universally quantified correspondence.

\textbf{Supplied mathematical context.} The source reports that harmonic weak Maass forms have a representation-theoretic interpretation through \((\mathfrak g,K)\)-modules, with finite-length representations whose composition factors are finite-dimensional representations or discrete series and their limits. It also reports that, in integral weight, iterated integrals connect with motivic geometry and can extend the notion of modular completion: mock modular forms and their generalizations are viewed as multivalued periods on a suitable moduli space of elliptic curves, and their completions as images under a single-valued projection. According to the supplied account, this formalism has been explicitly developed in the physics literature. After a minor weight adjustment using Maass operators, the completions are described there as linear combinations of products of essentially holomorphic and antiholomorphic functions, leading to tensor products of Harish--Chandra modules whose study is described as difficult. These are contextual claims of the supplied source, with their scope and hypotheses unspecified; they are not asserted here as unrestricted theorems.

The supplied original statement is truncated at ``polyharmonic weak Ma'' before its concluding remark. Definitions of higher-depth mixed mock modular forms, the relevant iterated integrals, the single-valued projection, and the intended correspondence are consequently not recoverable in full from this record. This draft preserves the research request and its qualifications without replacing the missing specifications by a conjectured theorem.
\subsection{Short English statement}
Find a precise connection between higher-depth mixed mock modular forms, iterated integrals, and polyharmonic weak Maass forms, and explain its representation-theoretic consequences. Begin by deciding whether iterated Eisenstein integrals correspond directly to polyharmonic Eisenstein series or need cusp-form or weak-form corrections. The source leaves the relevant classes, correspondence, and meaning of ``well-behaved'' unspecified, and its original text is truncated. The supplied open status is dated curation, not an independently verified conclusion.
\subsection{Sources}
\begin{itemize}
\item Supplied local UnsolvedMath registry record, permanent ID \(30006436\), title ``Higher-Depth Mock Modular Forms and Polyharmonic Maass Forms,'' source problem code \texttt{catalogue entry}, in the Oberwolfach Reports Open Problems collection. No individual proposer is named. The record gives proposed year \(2025\).
\item Supplied report citation: \emph{Cohomology Theories for Automorphic Forms and Enumerative Algebra} (2026), Oberwolfach Reports, pp. 2065--2124, DOI \texttt{10.4171/owr/2025/38}. Supplied link: \url{https://doi.org/10.4171/owr/2025/38}. The citation year is retained as supplied despite the different year in the DOI and proposed-year field.
\item Supplied literature source, ``Oberwolfach Report 38'': \url{https://ems.press/content/serial-article-files/52251}.
\item Supplied literature source, ``Matsusaka polyharmonic forms'': \url{https://doi.org/10.1007/s11139-018-0057-0}.
\item Supplied literature source, ``Alfes--Bringmann--Kudla classification'': \url{https://www.sciencedirect.com/science/article/pii/S0021869324004289}.
\item Literature assessment dated exactly \texttt{2026-08-22}: the supplied curation labels the problem open and describes the general connection through iterated integrals and its representation-theoretic explanation as unresolved, while noting substantial theories of polyharmonic weak Maass forms and higher-depth mock modular forms. This is supplied dated curation derived from the OWR record, not an independently verified literature assessment. Its time-sensitive status and cited claims require independent verification.
\item The registry labels the statement \texttt{corrected\_verified} and says only formula delimiters, notation, grammar, and typography were repaired while preserving mathematical content and scope. This supplied editorial label does not establish completeness or mathematical correctness. No source was visited in preparing this draft.
\end{itemize}
\end{document}
